\chapter*{Introduction}
\addcontentsline{toc}{chapter}{Introduction}

The nature of today’s world is greatly influenced by globalization,  
which not only connects economies and societies across continents, 
but also has a profound impact on cultures \cite{zizek1997}.
One of the main drivers of globalization is the transportation of goods. 
In order to transport goods effectively, containers are essential, 
as they allow for the easy movement of goods using various methods of transportation. 
They are used in trucks, ships, trains, and even airplanes, 
facilitating a smooth transition of goods when changing the mode of transport. 
The use of containers has tremendously reduced handling times and costs, making the global supply chain much more efficient \cite{bernhofen2016}.
Since the cost of shipping is typically calculated based on the number of containers, their volume, and their weight, 
any unused space consequently results in increased expenses. 
In large-scale logistics, where massive quantities of containers are shipped daily, 
even small improvements in packing efficiency can lead to significant cost savings as well as reduction in emissions. 
For example, in 2018, containers imported into the United States were, on average, loaded to only about 65\% of their capacity \cite{obyrne2021}.

In the context of computer science and optimization, container packing problems generally revolve around three key types \cite{martello2000}.
In 3D Bin Packing Problem, the goal is to pack a set of rectangular boxes into the smallest possible number of containers.
For Knapsack Loading Problem, each box has an associated value, and the objective is to select a subset of boxes that fits into a single container while maximizing the sum of their values.
The aim of Container Loading Problem is to pack all boxes into a single container, which has infinite length, while minimizing the total volume used.
All of these problems are indeed NP-hard \cite{martello2000}, so heuristic and approximation algorithms are usually used to find good solutions in a reasonable amount of time.

This thesis focuses on a variant of the 3D Bin Packing Problem in which rectangular boxes not only have given sizes,
but also an associated weight, and each container has a maximum load capacity. All orthogonal rotations of the boxes are allowed.
Furthermore, every box must be fully supported from below, either by the container floor or by the top surfaces of other boxes, 
ensuring structural stability of the packing. 
The primary objective is to minimize the number of containers used while respecting these constraints.

In this thesis, we build upon the genetic algorithm proposed by Gonçalves and Resende \cite{goncalves2013} and derive several 
algorithmic variants from it. In particular, we investigate whether its performance can be improved by 
incorporating mutation operators and memetic local search. Within this framework, 
we also examine the effect of replacing the box-packing order derived through evolution with several 
predefined ordering heuristics, as well as the effect of allowing evolution to choose the placement strategy 
for each box instead of using a single fixed strategy. In addition, these variants are compared with two other 
optimization approaches: mutation-based hill climbing and the physics-inspired method of simulated annealing \cite{eiben2016}. 
Finally, we investigate how the solutions change for different box density values under a fixed container weight constraint. 
The algorithms are evaluated according to solution quality, measured by fitness.

The thesis is structured as follows. 
Chapter~\ref{chap:problem} introduces the geometric terminology, formal definition of the packing problem and used datasets. 
Chapter~\ref{chap:packing} introduces a packing algorithm that deterministically constructs a complete solution from an individual and evaluates its fitness.
Chapter~\ref{chap:evolution} describes the evolutionary algorithms applied to the problem. 
Chapter~\ref{chap:experiments} reports on the experiments conducted and discusses the results. 
Finally, two appendices provide additional documentation: 
Appendix~\ref{app:user} contains user instructions, and Appendix~\ref{app:program} includes programming and implementation details.